% GENERATED FILE. Edit canonical lessons, then run npm run generate:books.
% canonical-source-hash: fnv1a64:1c9d7377222ca85c
% canonical-lessons: JA-C398-onsen, JA-C398-yuge, JA-C398-noboseru, JA-C398-tsukaru, JA-C398-kutsurogu

\chapter{The Hot Spring}
\label{ch:ja-the-hot-spring}
\hlchaptermodality{398}

\begin{chapteropening}
\textbf{By the end of this chapter:} \emph{I can enjoy a hot spring and say how it feels.}

Complete the last lesson of chapter 398: I can enjoy a hot spring and say how it feels.
\end{chapteropening}

\section[\texorpdfstring{\ja{おんせん} (onsen)}{おんせん (onsen)}]{\ja{おんせん} (onsen) — a hot spring}

\label{lesson:JA-C398-onsen}



\begin{quote}
Before the new one: say the Japanese for a day trip, then the Japanese for to look after.
\end{quote}

\subsection*{You'll want to know: \ja{おんせん}}
\textbf{\ja{おんせん}} — \emph{onsen} — \textquotedblleft{}a hot spring\textquotedblright{}.

\begin{quote}
\textbf{\ja{ふゆは} \ja{おんせんが} \ja{いちばん} \ja{いいです。}} — \emph{fuyu wa onsen ga ichiban ii desu} — \textquotedblleft{}In winter a hot spring is the best thing.\textquotedblright{}
\end{quote}

\begin{cousinweb}
\textbf{\ja{おんせん}} (\emph{onsen}) is Sino-Japanese (kango): two Chinese characters read with their on'yomi, readings borrowed from Chinese. They are \emph{on}, \textquotedblleft{}warm\textquotedblright{} (Middle Chinese \emph{'won}) and \emph{sen}, \textquotedblleft{}spring\textquotedblright{} (\emph{dzjwen}) (KANJIDIC2; Middle Chinese in Baxter's transcription).
\end{cousinweb}

\begin{grammarlens}[title={What you've built}]
One more word for getting around.
\end{grammarlens}

\subsection*{Your turn}
\begin{itemize}
\raggedright
  \item \emph{Say it:} \emph{onsen}
  \item \emph{Say it:} \emph{onsen}, once more
  \item \emph{Say it:} \emph{fuyu wa onsen ga ichiban ii desu} — In winter a hot spring is the best thing.
  \item \emph{Recall:} say the Japanese for a day trip, then the Japanese for to look after, then say \emph{onsen} again
\end{itemize}

\begin{tcolorbox}[breakable,colback=teal!4,colframe=teal!35!black,title={Before you move on}]
What does \ja{おんせん} mean? (\textquotedblleft{}A hot spring\textquotedblright{}.) Say it once more.
\end{tcolorbox}
\section[\texorpdfstring{\ja{ゆげ} (yuge)}{ゆげ (yuge)}]{\ja{ゆげ} (yuge) — rising steam; steam from hot water}

\label{lesson:JA-C398-yuge}



\begin{quote}
Before the new one: say the Japanese for to look after, then the Japanese for a hot spring.
\end{quote}

\subsection*{You'll want to know: \ja{ゆげ}}
\textbf{\ja{ゆげ}} — \emph{yuge} — \textquotedblleft{}rising steam; steam from hot water\textquotedblright{}.

\begin{quote}
\textbf{\ja{おんせんの} \ja{ゆげが} \ja{みえます。}} — \emph{onsen no yuge ga miemasu} — \textquotedblleft{}You can see the steam from the hot spring.\textquotedblright{}
\end{quote}

\begin{cousinweb}
\textbf{\ja{ゆげ}} (\emph{yuge}) is built from both layers of the vocabulary: the native Japanese (wago) \emph{yu}, \textquotedblleft{}hot water\textquotedblright{}, inherited rather than borrowed, and the Sino-Japanese \emph{ke}, \textquotedblleft{}air\textquotedblright{} (Middle Chinese \emph{khj+jH}), an on'yomi borrowed from Chinese, joined as \emph{yuge} (JMdict; KANJIDIC2; Middle Chinese in Baxter's transcription).
\end{cousinweb}

\begin{grammarlens}[title={What you've built}]
One more word for getting around.
\end{grammarlens}

\subsection*{Your turn}
\begin{itemize}
\raggedright
  \item \emph{Say it:} \emph{yuge}
  \item \emph{Say it:} \emph{yuge}, once more
  \item \emph{Say it:} \emph{onsen no yuge ga miemasu} — You can see the steam from the hot spring.
  \item \emph{Recall:} say the Japanese for to look after, then the Japanese for a hot spring, then say \emph{yuge} again
\end{itemize}

\begin{tcolorbox}[breakable,colback=teal!4,colframe=teal!35!black,title={Before you move on}]
What does \ja{ゆげ} mean? (\textquotedblleft{}Rising steam; steam from hot water\textquotedblright{}.) Say it once more.
\end{tcolorbox}
\section[\texorpdfstring{\ja{のぼせる} (noboseru)}{のぼせる (noboseru)}]{\ja{のぼせる} (noboseru) — to feel dizzy (from the heat of a bath)}

\label{lesson:JA-C398-noboseru}



\begin{quote}
Before the new one: say the Japanese for a hot spring, then the Japanese for rising steam.
\end{quote}

\subsection*{You'll want to know: \ja{のぼせる}}
\textbf{\ja{のぼせる}} — \emph{noboseru} — \textquotedblleft{}to feel dizzy (from the heat of a bath)\textquotedblright{}.

\begin{quote}
\textbf{\ja{ながく} \ja{おふろに} \ja{はいって、のぼせました。}} — \emph{nagaku ofuro ni haitte, nobosemashita} — \textquotedblleft{}I stayed in the bath too long and felt dizzy.\textquotedblright{}
\end{quote}

Polite \textbf{\ja{のぼせます}}, \emph{nobosemasu}; an ichidan verb. Past (polite): \textbf{\ja{あにが} \ja{おふろで} \ja{のぼせました。}}, \emph{ani ga ofuro de nobosemashita}.

\begin{cousinweb}
\textbf{\ja{のぼせる}} (\emph{noboseru}) is a native Japanese word (wago), inherited rather than borrowed from Chinese (JMdict).
\end{cousinweb}

\begin{grammarlens}[title={What you've built}]
One more word for getting around.
\end{grammarlens}

\subsection*{Your turn}
\begin{itemize}
\raggedright
  \item \emph{Say it:} \emph{noboseru}
  \item \emph{Say it:} \emph{noboseru}, once more
  \item \emph{Say it:} \emph{nagaku ofuro ni haitte, nobosemashita} — I stayed in the bath too long and felt dizzy.
  \item \emph{Recall:} say the Japanese for a hot spring, then the Japanese for rising steam, then say \emph{noboseru} again
\end{itemize}

\begin{tcolorbox}[breakable,colback=teal!4,colframe=teal!35!black,title={Before you move on}]
What does \ja{のぼせる} mean? (\textquotedblleft{}To feel dizzy (from the heat of a bath)\textquotedblright{}.) Say it once more.
\end{tcolorbox}
\section[\texorpdfstring{\ja{つかる} (tsukaru)}{つかる (tsukaru)}]{\ja{つかる} (tsukaru) — to soak (in water)}

\label{lesson:JA-C398-tsukaru}



\begin{quote}
Before the new one: say the Japanese for rising steam, then the Japanese for to feel dizzy.
\end{quote}

\subsection*{You'll want to know: \ja{つかる}}
\textbf{\ja{つかる}} — \emph{tsukaru} — \textquotedblleft{}to soak (in water)\textquotedblright{}.

\begin{quote}
\textbf{\ja{あつい} \ja{ゆに} \ja{ゆっくり} \ja{つかりました。}} — \emph{atsui yu ni yukkuri tsukarimashita} — \textquotedblleft{}I had a long soak in the hot water.\textquotedblright{}
\end{quote}

Polite \textbf{\ja{つかります}}, \emph{tsukarimasu}; a godan verb. Past (polite): \textbf{\ja{ちちが} \ja{おんせんに} \ja{つかりました。}}, \emph{chichi ga onsen ni tsukarimashita}.

\begin{cousinweb}
\textbf{\ja{つかる}} (\emph{tsukaru}) is a native Japanese word (wago), inherited rather than borrowed from Chinese; its kanji was attached later as a kun'yomi (JMdict).
\end{cousinweb}

\begin{grammarlens}[title={What you've built}]
One more word for getting around.
\end{grammarlens}

\subsection*{Your turn}
\begin{itemize}
\raggedright
  \item \emph{Say it:} \emph{tsukaru}
  \item \emph{Say it:} \emph{tsukaru}, once more
  \item \emph{Say it:} \emph{atsui yu ni yukkuri tsukarimashita} — I had a long soak in the hot water.
  \item \emph{Recall:} say the Japanese for rising steam, then the Japanese for to feel dizzy, then say \emph{tsukaru} again
\end{itemize}

\begin{tcolorbox}[breakable,colback=teal!4,colframe=teal!35!black,title={Before you move on}]
What does \ja{つかる} mean? (\textquotedblleft{}To soak (in water)\textquotedblright{}.) Say it once more.
\end{tcolorbox}
\section[\texorpdfstring{\ja{くつろぐ} (kutsurogu)}{くつろぐ (kutsurogu)}]{\ja{くつろぐ} (kutsurogu) — to relax; to make oneself at home}

\label{lesson:JA-C398-kutsurogu}



\begin{quote}
Before the new one: say the Japanese for to feel dizzy, then the Japanese for to soak.
\end{quote}

\subsection*{You'll want to know: \ja{くつろぐ}}
\textbf{\ja{くつろぐ}} — \emph{kutsurogu} — \textquotedblleft{}to relax; to make oneself at home\textquotedblright{}.

\begin{quote}
\textbf{\ja{りょかんの} \ja{へやで} \ja{くつろぎました。}} — \emph{ryokan no heya de kutsurogimashita} — \textquotedblleft{}We relaxed in our room at the inn.\textquotedblright{}
\end{quote}

Polite \textbf{\ja{くつろぎます}}, \emph{kutsurogimasu}; a godan verb. Past (polite): \textbf{\ja{かぞくが} \ja{へやで} \ja{くつろぎました。}}, \emph{kazoku ga heya de kutsurogimashita}.

\begin{cousinweb}
\textbf{\ja{くつろぐ}} (\emph{kutsurogu}) is a native Japanese word (wago), inherited rather than borrowed from Chinese; its kanji was attached later as a kun'yomi (JMdict).
\end{cousinweb}

\begin{grammarlens}[title={What you've built}]
One more word for getting around.
\end{grammarlens}

\subsection*{Your turn}
\begin{itemize}
\raggedright
  \item \emph{Say it:} \emph{kutsurogu}
  \item \emph{Say it:} \emph{kutsurogu}, once more
  \item \emph{Say it:} \emph{ryokan no heya de kutsurogimashita} — We relaxed in our room at the inn.
  \item \emph{Recall:} say the Japanese for to feel dizzy, then the Japanese for to soak, then say \emph{kutsurogu} again
\end{itemize}

\begin{tcolorbox}[breakable,colback=teal!4,colframe=teal!35!black,title={Before you move on}]
What does \ja{くつろぐ} mean? (\textquotedblleft{}To relax; to make oneself at home\textquotedblright{}.) Say it once more.
\end{tcolorbox}
